\chapter{Modèle économique et stratégie de rentabilisation}
\label{ch:economie}

LexHub est un SaaS facturable. Ce chapitre décrit ce qui est vendu, comment l'essai fonctionne, comment recharger des crédits, comment le client est informé de son usage, et qui a le droit de payer. Les tarifs publics sont ceux de l'API \texttt{/api/billing/plans} (production : \url{https://lexhub.space}).

\section{Plans et essai hybride}

\begin{table}[H]
\centering
\caption{Plans LexHub (production)}
\label{tab:plans}
\small
\begin{tabular}{@{}lrrrl@{}}
\toprule
\textbf{Plan} & \textbf{MAD/mois} & \textbf{MAD/an} & \textbf{Essai} & \textbf{Plaidoiries / crédits} \\
\midrule
Découverte & 0 & 0 & --- & 3 / 50 \\
Solo & 299 & 2\,990 & 14~j & 35 / 600 \\
Cabinet & 699 & 6\,990 & 30~j & 120 / 2\,000 \\
Entreprise & 1\,290 & 12\,900 & 30~j & 350 / 5\,000 \\
\bottomrule
\end{tabular}
\end{table}

L'\textbf{essai hybride} ne facture pas ($0$~MAD) mais n'ouvre \textbf{pas} le forfait plein. Les plafonds d'essai sont inférieurs aux quotas payants : Solo $400$ crédits / $25$ plaidoiries ; Cabinet $1\,000$ / $60$ ; Entreprise $1\,800$ / $120$. Les \textbf{packs} ne sont pas proposés pendant l'essai (pas de vente de recharge). En revanche, si le forfait d'essai est épuisé, le client peut \textbf{souscrire} via les tarifs pour continuer immédiatement, ou attendre la date de fin d'essai.

\section{Repère IA du jour (tous les plans)}
\label{sec:repere-jour}

Indépendamment de l'essai ou de l'abonnement payant, chaque cabinet dispose d'un \textbf{repère de crédits IA par jour} (anti-abus). Ce n'est \textbf{pas} un impayé et ce n'est \textbf{pas} la fin de l'essai : c'est une information d'usage. La même logique s'applique à Découverte, Solo, Cabinet et Entreprise.

\begin{table}[H]
\centering
\caption{Repère IA du jour (essai vs abonnement payé)}
\label{tab:daily-cap}
\small
\begin{tabular}{@{}lrr@{}}
\toprule
\textbf{Plan} & \textbf{Essai} & \textbf{Payé} \\
\midrule
Découverte & --- & 12 \\
Solo & 35 & 50 \\
Cabinet & 90 & 180 \\
Entreprise & 105 & 300 \\
\bottomrule
\end{tabular}
\end{table}

Lorsque l'action suivante ferait dépasser ce repère, une popup indique clairement le contexte (\textit{essai} ou \textit{abonnement}, nom du plan, compteur du type $33/35$) et propose deux choix : \textbf{Continuer maintenant} (l'IA part ; les crédits du mois ou de l'essai restent débités) ou \textbf{Attendre demain} (compteur du jour remis à zéro). Après « Continuer », la session se souvient pour le reste de la journée. Une pastille UI (crédits restants $\cdot$ aujourd'hui) et la barre \emph{IA aujourd'hui} dans Paramètres $\rightarrow$ Plan \& usage complètent l'information.

\section{Packs, sièges et usage à la demande}

Trois objets restent distincts. L'\textbf{abonnement} donne l'accès et le forfait du mois. Les \textbf{packs} (500 crédits / $110$~MAD, $1\,000$ / $200$~MAD, $2\,500$ / $450$~MAD) ne s'achètent que si l'abonnement est \textbf{actif et payant} ; chaque lot expire au bout de $12$~mois (consommation FIFO). L'\textbf{usage à la demande} (PAYG) est désactivé par défaut ($\approx 0{,}25$~MAD par crédit) et indisponible sur Découverte.

Ordre de consommation : forfait du mois (ou quotas d'essai), puis packs, puis PAYG s'il est activé ; sinon popup de résolution (section suivante).

Sur le plan Cabinet, $0$ à $5$ sièges avocats supplémentaires coûtent $99$~MAD par siège et par mois.

Le paiement carte marocain (CMI / PayZone) dispose d'un \textbf{environnement technique prêt}. En production actuelle, le parcours de paiement reste une \textbf{simulation} (activation opérateur / identifiant marchand non finalisée) : aucun débit réel n'est présenté comme déjà en service.

\section{Quand le forfait est épuisé --- popups de guidance}
\label{sec:forfait-epuise}

Deux situations se distinguent du repère du jour. Ici le forfait (essai ou mois) est vraiment terminé : l'action IA est refusée tant que le client n'a pas choisi une issue.

\begin{table}[H]
\centering
\caption{Guidance client lorsque le forfait est épuisé}
\label{tab:quota-guidance}
\small
\renewcommand{\arraystretch}{1.15}
\begin{tabular}{@{}p{3.4cm}p{4.4cm}p{4.6cm}@{}}
\toprule
\textbf{Situation} & \textbf{Popup} & \textbf{Choix proposés} \\
\midrule
Forfait d'essai épuisé &
« Forfait d'essai épuisé » + date de fin &
Attendre la fin d'essai (pas de packs) \textbf{ou} souscrire (voir les tarifs) \\
Forfait du mois (abonné) épuisé &
« Forfait du mois épuisé » &
(1) Activer le PAYG \newline
(2) Acheter un pack \newline
(3) Passer au plan supérieur \newline
sinon attendre le mois suivant \\
Découverte épuisée &
Blocage quota &
Upgrade Solo / Cabinet (pas de packs, pas de PAYG) \\
\bottomrule
\end{tabular}
\end{table}

Des toasts (solde bas / crédits épuisés) et la bannière d'abonnement (essai bientôt fini, renouvellement, \texttt{past\_due}) renforcent cette guidance sans remplacer les popups.

\section{Parrainage cabinet $\rightarrow$ cabinet}
\label{sec:parrainage}

Chaque cabinet a un code (\texttt{LH-\ldots}) et un lien \texttt{?ref=}. Le filleul reçoit $14$ jours d'essai supplémentaires au démarrage Solo ou Cabinet. Le parrain reçoit $30$ jours lorsque l'abonnement payant du filleul commence après l'essai. L'auto-parrainage est interdit (plafond : $12$ filleuls par an). Ce mécanisme n'est pas une invitation d'équipe : il ne donne aucun accès aux dossiers du parrain.

\section{Répartition payeur / consommateur}

Seul l'\textbf{admin cabinet} change de plan, paie, achète des packs, active le PAYG ou ajoute des sièges. Les avocats et les assistants consomment le forfait et voient les quotas en lecture seule, sans débiter une carte. Une résiliation prend effet à la fin de la période déjà payée.
